\chapter{De ontologie van het kunstwerk en zijn hermeneutische betekenis}

\section{Het spel als leidraad voor de ontologische verklaring}

\subsection{Het begrip spel}

Als uitgangspunt kies ik een idee dat een belangrijke rol heeft gespeeld in de esthetiek: het begrip spel. Ik wil dit begrip bevrijden van de subjectieve betekenis die het bij Kant en Schiller heeft, en die de hele moderne esthetiek en filosofie van de mens domineert. Wanneer we over spel spreken in relatie tot de ervaring van kunst, bedoelen we daarmee noch de gerichtheid noch zelfs de gemoedstoestand van de schepper of van hen die genieten van het kunstwerk, noch de vrijheid van een subjectiviteit die zich in spel begeeft, maar de wijze van zijn van het kunstwerk zelf. Bij de analyse van het esthetische bewustzijn hebben we al opgemerkt dat het opvatten van esthetisch bewustzijn als iets dat een object tegenover zich heeft, niet recht doet aan de werkelijke situatie. Om die reden is het begrip spel van belang in mijn uiteenzetting.

We kunnen zeker een onderscheid maken tussen spel en het gedrag van de speler, dat als zodanig behoort tot de andere vormen van subjectief gedrag. Zo kan men zeggen dat spel voor de speler niet serieus is: daarom speelt hij. We kunnen proberen het begrip spel vanaf dit standpunt te definiëren. Wat louter spel is, is niet serieus. Spel heeft een bijzondere relatie tot het serieuze. Het is niet alleen zo dat het serieuze er zijn ``doel'' geeft: we spelen ``ter ontspanning'', zoals Aristoteles zegt. Belangrijker is dat spel zelf zijn eigen, zelfs heilige, serieuziteit in zich draagt. Toch zijn in het spel al die doelgerichte relaties die het actieve en bezorgde bestaan bepalen niet eenvoudigweg verdwenen, maar op merkwaardige wijze opgeschort. De speler beseft dat spel slechts spel is en plaatsvindt binnen een wereld die wordt gekenmerkt door de ernst van doelgerichte bezigheden. Maar hij weet dit niet op zo'n manier dat hij als speler deze relatie tot het serieuze daadwerkelijk bedoelt. Spel vervult zijn doel alleen als de speler zich in het spel verliest. Ernst is niet slechts iets dat ons van het spel afleidt; integendeel, ernst in het spel is noodzakelijk om het spel volledig spel te laten zijn. Iemand die het spel niet serieus neemt, is een spelbreker. De wijze van zijn van het spel staat niet toe dat de speler zich tegenover het spel gedraagt alsof het een object is. De speler weet heel goed wat spel is, en dat wat hij doet ``slechts een spel'' is; maar hij weet niet precies wat hij ``weet'' als hij dat weet.

Onze vraag naar de aard van het spel zelf kan daarom geen antwoord vinden als we deze zoeken in de subjectieve reflectie van de speler. In plaats daarvan onderzoeken we de wijze van zijn van het spel als zodanig. We hebben gezien dat het niet het esthetische bewustzijn is, maar de ervaring van kunst — en dus de vraag naar de wijze van zijn van het kunstwerk — die het voorwerp van ons onderzoek moet zijn. Maar dit was precies de ervaring van het kunstwerk die ik in tegenstelling stelde tot het nivelleringsproces van het esthetische bewustzijn: namelijk dat het kunstwerk geen object is dat voor een subject op zichzelf staat. Integendeel, het kunstwerk heeft zijn ware zijn in het feit dat het een ervaring wordt die de persoon die het ervaart, verandert. Het ``subject'' van de ervaring van kunst, dat wat blijft en standhoudt, is niet de subjectiviteit van de persoon die het ervaart, maar het werk zelf. Hier wordt de wijze van zijn van het spel relevant. Want spel heeft zijn eigen wezen, onafhankelijk van het bewustzijn van hen die spelen. Spel — inderdaad, spel in eigenlijke zin — bestaat ook wanneer de thematische horizon niet wordt beperkt door enig zijn-voor-zich van subjectiviteit, en waar er geen subjecten zijn die zich ``speels'' gedragen.

De spelers zijn niet de subjecten van het spel; in plaats daarvan komt het spel slechts tot voorstelling door de spelers. Dit blijkt al uit het gebruik van het woord, met name uit de vele metaforische toepassingen, die Buytendijk in het bijzonder heeft opgemerkt.

Hier, zoals altijd, heeft het metaforische gebruik methodologische voorrang. Wanneer een woord wordt toegepast op een sfeer waartoe het oorspronkelijk niet behoorde, komt de werkelijke ``originele'' betekenis duidelijk naar voren. De taal heeft vooraf de abstractie verricht die als zodanig de taak is van de begripsmatige analyse. Het denken hoeft deze voorafgaande prestatie slechts te benutten.

Hetzelfde geldt ook voor etymologieën. Deze zijn weliswaar veel minder betrouwbaar, omdat het abstracties zijn die niet door de taal zelf, maar door de taalkunde zijn bereikt, en nooit volledig kunnen worden geverifieerd door de taal zelf, dat wil zeggen door het daadwerkelijke gebruik. Daarom zijn etymologieën, zelfs wanneer ze correct zijn, geen bewijzen, maar voorbereidende prestaties voor de begripsmatige analyse, en verkrijgen ze pas een vaste grondslag in een dergelijke analyse.

Als we onderzoeken hoe het woord ``spel'' wordt gebruikt en ons concentreren op de zogenaamde metaforische betekenissen, dan vinden we uitdrukkingen als het spel van het licht, het spel van de golven, het spel van tandwielen of onderdelen van een machine, het samenspel van ledematen, het spel van krachten, het spel van muggen, zelfs een woordspel. In elk geval wordt bedoeld een heen-en-weer beweging die niet gebonden is aan enig doel dat haar zou beëindigen. Overigens betekende het woord ``Spiel'' oorspronkelijk ``dans'', en komt het nog steeds voor in veel woordvormen (bijvoorbeeld in ``Spielmann'', ``jongleur''). De beweging van het spelen heeft geen doel dat haar beëindigt; in plaats daarvan vernieuwt zij zich in constante herhaling. De heen-en-weer beweging is kennelijk zo centraal voor de definitie van spel dat het niet uitmaakt wie of wat deze beweging uitvoert. De beweging van het spel als zodanig heeft, zo zou men kunnen zeggen, geen substraat. Het is het spel dat gespeeld wordt — het is onbelangrijk of er al dan niet een subject is dat het speelt. Het spel is het gebeuren van de beweging als zodanig. Zo spreken we van het spel van kleuren en bedoelen daarmee niet alleen dat de ene kleur tegen de andere speelt, maar dat er één proces of gezicht is dat een wisselende verscheidenheid aan kleuren toont.

De zijnswijze van het spel is dus niet zodanig dat er, opdat er gespeeld kan worden, eerst een subject moet zijn dat zich spelenderwijs gedraagt. Integendeel, de oorspronkelijke betekenis van ‘spelen’ is de mediale betekenis. Daarom zeggen we dat er ergens of op een bepaald moment ‘iets speelt’, dat iets ‘in het spel is’ of dat zich iets ‘afspeelt’.

Deze taalkundige observatie lijkt mij een indirecte aanwijzing dat spel niet moet worden opgevat als iets wat een persoon doet. Voor zover het de taal betreft, is het werkelijke subject van het spel kennelijk niet de subjectiviteit van een individu dat, naast andere activiteiten, ook speelt, maar het spel zelf. Maar we zijn zo gewend om verschijnselen als spelen te relateren aan de sfeer van de subjectiviteit en de wijze waarop deze handelt, dat we gesloten blijven voor deze aanwijzingen uit de geest van de taal.

Toch heeft modern antropologisch onderzoek de aard van het spel zo breed opgevat dat het bijna voorbij is aan het opvatten van spel als subjectiviteit. Huizinga heeft het element van spel in alle culturen onderzocht en heeft met name de verbinding tussen kinderspel, dierenspel en ``heilig spel'' uitgewerkt. Dat leidde hem ertoe om de merkwaardige onbeslistheid van het spelende bewustzijn te herkennen, die het absoluut onmogelijk maakt om te beslissen tussen geloof en niet-geloof. ``De wildeman zelf kent geen begripsmatig onderscheid tussen zijn en spelen; hij weet niets van identiteit, van afbeelding of symbool. En daarom kan men vragen of de geestestoestand van de wildeman bij zijn heilige plechtigheden niet het best wordt begrepen door spel als de primaire term te behouden. In ons begrip van spel is het verschil tussen geloof en veinzen opgelost.''

Hier wordt de \textit{voorrang van het spel boven het bewustzijn van de speler} fundamenteel erkend, en inderdaad worden ook de ervaringen van spel die psychologen en antropologen beschrijven op nieuwe wijze belicht als men uitgaat van de mediale betekenis van het woord ``spelen''. Spel vertegenwoordigt kennelijk een orde waarin de heen-en-weer beweging van het spel vanzelf volgt. Het hoort bij het spel dat de beweging niet alleen zonder doel of bedoeling is, maar ook zonder inspanning. Het gebeurt, zo zou men kunnen zeggen, vanzelf. De lichtheid van het spel — die natuurlijk niet betekent dat er sprake is van een werkelijke afwezigheid van inspanning, maar fenomelogisch slechts verwijst naar de afwezigheid van spanning — wordt subjectief ervaren als ontspanning. De structuur van het spel absorbeert de speler in zichzelf en bevrijdt hem daardoor van de last om het initiatief te nemen, wat de werkelijke spanning van het bestaan uitmaakt. Dit zien we ook in de spontane neiging tot herhaling die bij de speler ontstaat en in de constante zelfvernieuwing van het spel, die zijn vorm beïnvloedt (bijvoorbeeld het refrein).

Het feit dat de wijze van zijn van het spel zo dicht staat bij de mobiele vorm van de natuur, stelt ons in staat een belangrijke methodologische conclusie te trekken. Het is kennelijk niet correct om te zeggen dat dieren \textit{ook} spelen, noch om te zeggen dat, figuurlijk gesproken, water en licht \textit{eveneens} spelen. Integendeel, we kunnen zeggen dat \textit{de mens} ook speelt. Zijn spelen is eveneens een natuurlijk proces. De betekenis van zijn spel is, precies omdat — en voor zover — hij deel uitmaakt van de natuur, een zuivere zelfpresentatie. Zo wordt in deze sfeer het uiteindelijk zinloos om een onderscheid te maken tussen letterlijke en metaforische toepassing.

Maar het belangrijkste is dat het zijn van het kunstwerk verbonden is met de mediale betekenis van spel (\textit{Spiel}: ook spelletje en drama). Voor zover de natuur zonder doel en bedoeling is, net zoals zij zonder inspanning is, is zij een constant zich vernieuwend spel, en kan zij daarom als model voor de kunst verschijnen. Zo schrijft Friedrich Schlegel: ``Alle heilige spellen van de kunst zijn slechts verre nabootsingen van het oneindige spel van de wereld, het eeuwig zichzelf scheppende werk van de kunst.''

Een andere vraag die Huizinga bespreekt, wordt ook verduidelijkt door de fundamentele rol van de heen-en-weer beweging van het spel: namelijk het speelse karakter van de wedstrijd. Het is waar dat de deelnemer aan de wedstrijd zichzelf niet als spelend beschouwt. Maar door de wedstrijd ontstaat de gespannen heen-en-weer beweging waaruit de overwinnaar tevoorschijn komt, en zo wordt het geheel een spel. De heen-en-weer beweging hoort kennelijk zo wezenlijk bij het spel dat er in laatste instantie geen spel mogelijk is als men alleen speelt. Opdat er een spel kan zijn, moet er altijd iets zijn, niet noodzakelijkerwijs letterlijk een andere speler, maar iets anders waartegen de speler speelt en dat automatisch reageert op zijn zet met een tegenzet. Zo kiest de kat die speelt voor de wolbal omdat deze meespeelt, en zullen balspellen voor altijd bij ons blijven omdat de bal vrij beweeglijk is in alle richtingen en ogenschijnlijk verrassende dingen uit eigen beweging doet.

In gevallen waarin het de menselijke subjectiviteit is die speelt, ervaart de speler zelf op een bijzondere manier de voorrang van het spel boven de spelers die eraan deelnemen. Opnieuw zijn het de onjuiste, metaforische toepassingen van het woord die de meeste informatie bieden over zijn eigenlijke wezen. Zo zeggen we van iemand dat hij speelt met mogelijkheden of met plannen. Wat we daarmee bedoelen, is duidelijk. Hij heeft zich nog niet verbonden met de mogelijkheden als met serieuze doelen. Hij heeft nog de vrijheid om voor de ene of de andere mogelijkheid te kiezen. Aan de andere kant is deze vrijheid niet zonder gevaar. Integendeel, het spel zelf is een risico voor de speler. Men kan alleen spelen met serieuze mogelijkheden. Dit betekent kennelijk dat men zo in hen op kan gaan dat zij, zo zou men kunnen zeggen, hem overspelen en de overhand over hem krijgen. De aantrekkingskracht die het spel op de speler uitoefent, ligt in dit risico. Men geniet van een vrijheid van keuze die tegelijkertijd bedreigd is en onherroepelijk beperkt. Men hoeft alleen maar te denken aan legpuzzels, geduldspellen, en dergelijke. Maar hetzelfde geldt voor serieuze zaken. Als iemand, omwille van het genot van zijn eigen vrijheid van keuze, het nemen van dringende beslissingen vermijdt of speelt met mogelijkheden die hij niet serieus overweegt en die daarom geen risico inhouden dat hij voor hen zal kiezen en zich daardoor beperken, dan zeggen we dat hij slechts ``met het leven speelt'' (verspeelt).

Dit suggereert een algemene eigenschap van de aard van het spel die zich weerspiegelt in het spelen: al het spelen is een gespeeld worden. De aantrekkingskracht van een spel, de fascinatie die het uitoefent, bestaat juist hierin dat het spel de spelers beheerst. Zelfs in het geval van spellen waarin men probeert taken uit te voeren die men zichzelf heeft opgelegd, is er een risico dat deze niet ``werken'', ``lukken'' of ``opnieuw lukken'', en dat is juist de aantrekkingskracht van het spel. Wie ``probeert'', is in feite degene die op de proef wordt gesteld. Het ware subject van het spel (dit blijkt juist uit die ervaringen waarin er maar één speler is) is niet de speler, maar het spel zelf. Wat de speler in zijn ban houdt, hem in het spel trekt en daar houdt, is het spel zelf.

Dit blijkt ook uit het feit dat elk spel zijn eigen, specifieke geest heeft. Maar zelfs dit verwijst niet naar de gemoedstoestand of de geestelijke toestand van hen die het spel spelen. Integendeel, de verscheidenheid aan geestelijke houdingen die men waarneemt bij het spelen van verschillende spellen, en in de wens om ze te spelen, is het gevolg en niet de oorzaak van de verschillen tussen de spellen zelf. Spellen verschillen van elkaar in hun geest. De reden hiervoor is dat de heen-en-weer beweging die het spel vormgeeft, op verschillende manieren is gestructureerd. De bijzondere aard van een spel ligt in de regels en voorschriften die voorschrijven op welke manier het speelveld wordt ingevuld. Dit geldt universeel, waar ook een spel wordt gespeeld. Het is bijvoorbeeld waar voor het spel van fonteinen en van spelende dieren. Het speelveld waarop het spel wordt gespeeld, is zo te zeggen vastgesteld door de aard van het spel zelf en wordt veel meer bepaald door de structuur die de beweging van het spel van binnenuit regelt dan door wat het van buitenaf tegenkomt — dat wil zeggen, de grenzen van de open ruimte — die de beweging van buitenaf beperken.

Naast deze algemene bepalende factoren, lijkt het mij kenmerkend voor menselijk spel dat het iets speelt. Dat betekent dat de structuur van beweging waaraan het zich onderwerpt, een bepaalde kwaliteit heeft die de speler ``kiest''. Ten eerste scheidt hij zijn speelgedrag uitdrukkelijk van zijn ander gedrag door te willen spelen. Maar zelfs binnen zijn bereidheid om te spelen, maakt hij een keuze. Hij kiest dit spel in plaats van dat. Daardoor is de ruimte waarin de beweging van het spel plaatsvindt niet eenvoudigweg de open ruimte waarin men ``zich uitspeelt'', maar een ruimte die speciaal is afgebakend en gereserveerd voor de beweging van het spel. Menselijk spel vereist een speelveld. Het afbakenen van het speelveld — net zoals het afbakenen van heilige terreinen, zoals Huizinga terecht opmerkt — scheidt de sfeer van het spel af als een gesloten wereld, een wereld zonder overgang en bemiddeling naar de wereld van doelen.

Dat al het spelen het spelen van iets is, geldt hier, waar de geordende heen-en-weer beweging van het spel wordt bepaald als één soort gedrag (\textit{Verhalten}) onder andere. Een mens die speelt, is, zelfs in zijn spel, nog steeds iemand die zich gedraagt, ook al bestaat het eigenlijke wezen van het spel erin dat hij zich ontlast van de spanning die hij voelt in zijn doelgerichte gedrag. Dit bepaalt nauwkeuriger waarom spelen altijd een spelen van iets is. Elk spel stelt de mens die het speelt voor een taak. Hij kan niet genieten van de vrijheid om zich uit te spelen zonder de doelen van zijn doelgerichte gedrag om te zetten in louter taken van het spel. Zo geeft het kind zichzelf een taak in het spelen met een bal, en dergelijke taken zijn speels omdat het doel van het spel niet echt het oplossen van de taak is, maar het ordenen en vormgeven van de beweging van het spel zelf.

Kennelijk hangt de kenmerkende lichtheid en het gevoel van verlichting die we in het spel vinden af van de specifieke aard van de taak die door het spel wordt gesteld en komt voort uit het oplossen daarvan.

Men kan zeggen dat het succesvol uitvoeren van een taak deze ``voorstelt'' (\textit{stellt sie dar}). Deze formulering dringt zich met name op in het geval van een spel, want hier betekent het volbrengen van de taak niet dat deze naar een doelgerichte context verwijst. Spel is inderdaad beperkt tot zichzelf voorstellen. Zo is zijn wijze van zijn zelfpresentatie. Maar zelfpresentatie is een universele ontologische eigenschap van de natuur. We weten vandaag de dag hoe onvoldoende opvattingen over biologische doeleinden zijn als het gaat om het begrijpen van de vorm van levende wezens. Zo is het ook een onvoldoende benadering om te vragen wat de levensfunctie en het biologische doel van spel is. Allereerst is spel zelfpresentatie.

Zoals we hebben gezien, hangt de zelfpresentatie van menselijk spel af van het feit dat het gedrag van de speler verbonden is met de schijnbare doelen van het spel, maar de ``betekenis'' van deze doelen hangt in feite niet af van het feit dat ze worden bereikt. Integendeel, door zich in te spannen voor de taak van het spel, speelt men zich in feite uit. De zelfpresentatie van het spel omvat dat de speler, zo zou men kunnen zeggen, zijn eigen zelfpresentatie bereikt door iets te spelen — dat wil zeggen, voor te stellen. Alleen omdat spel altijd presentatie is, kan menselijk spel de voorstelling zelf tot taak van een spel maken. Zo zijn er spellen die men voorstellingsspellen kan noemen, ofwel omdat zij, door hun gebruik van betekenisvolle toespelingen, iets van voorstelling in zich dragen (bijvoorbeeld ``Smid, Kleermaker, Soldaat, Zeeman''), ofwel omdat het spel zelf bestaat uit het voorstellen van iets (bijvoorbeeld wanneer kinderen auto's spelen).

Alle presentatie is potentieel een voorstelling voor iemand. Dat deze mogelijkheid bedoeld is, is het kenmerkende van kunst als spel. De gesloten wereld van het spel laat zo te zeggen een van zijn muren zakken. Een religieuze ritus en een toneelspel stellen kennelijk niet op dezelfde manier voor als een kind dat speelt. Hun zijn is niet uitgeput door het feit dat zij zichzelf presenteren, want tegelijkertijd wijzen zij voorbij zichzelf naar het publiek dat deelnemt door toe te kijken. Spel is hier niet langer louter de zelfpresentatie van een geordende beweging, noch louter voorstelling waarin het spelende kind volledig opgaat, maar het is ``voorstellen voor iemand''. De gerichtheid die eigen is aan alle voorstelling komt hier op de voorgrond en is constitutief voor het zijn van kunst.

Over het algemeen, hoezeer spellen in wezen voorstellingen zijn en hoezeer de spelers zichzelf daarin voorstellen, spellen zijn niet voor iemand bedoeld — dat wil zeggen, ze zijn niet gericht op een publiek. Kinderen spelen voor zichzelf, zelfs wanneer zij iets voorstellen. En zelfs die spellen (bijvoorbeeld sporten) die voor toeschouwers worden gespeeld, zijn niet op hen gericht. Integendeel, wedstrijden lopen het risico hun echte speelse karakter juist te verliezen door spectaculair te worden. Een processie als onderdeel van een religieuze ritus is meer dan een schouwspel, omdat de werkelijke betekenis ervan is om de hele religieuze gemeenschap te omvatten. En toch is een religieuze handeling een echte voorstelling voor de gemeenschap; en evenzo is een drama een soort spelen dat van nature om een publiek vraagt. De presentatie van een god in een religieuze ritus, de presentatie van een mythe in een toneelstuk, zijn spel niet alleen in die zin dat de deelnemende spelers volledig opgaan in het presentatiespel en daarin hun verhoogde zelfpresentatie vinden, maar ook in die zin dat de spelers een betekenisvol geheel voor een publiek voorstellen. Zo is het niet echt de afwezigheid van een vierde wand die het toneelspel in een voorstelling verandert. Integendeel, openheid naar de toeschouwer toe maakt deel uit van de geslotenheid van het spel. Het publiek voltooit slechts wat het spel als zodanig is.

Dit punt laat zien hoe belangrijk het is om spel te begrijpen als een proces dat zich ``tussen'’ de deelnemers voltrekt. We hebben gezien dat spel zijn bestaan niet heeft in het bewustzijn of de houding van de speler. Integendeel, het spel trekt de speler in zijn eigen sfeer en vervult hem met zijn geest. De speler ervaart het spel als een werkelijkheid die hem overstijgt. Dat geldt des te sterker wanneer het spel zelf bedoeld is als zo een werkelijkheid — bijvoorbeeld bij toneelspel, dat zich als \textit{opvoering aan een publiek} presenteert.

Zelfs een toneelspel blijft een spel — dat wil zeggen, het heeft de structuur van een spel, die is die van een gesloten wereld. Maar hoezeer een religieus of profaan spel ook een wereld voorstelt die volledig in zichzelf gesloten is, het is alsof het openstaat naar de toeschouwer, in wie het zijn volledige betekenis bereikt. De spelers spelen hun rollen zoals in elk spel, en zo wordt het spel voorgesteld, maar het spel zelf is het geheel, bestaande uit spelers en toeschouwers. In feite wordt het op de juiste manier ervaren door, en stelt het zich voor (zoals het ``bedoeld'' is) aan iemand die niet in het spel meespeelt, maar het bekijkt. In hem wordt het spel, zo zou men kunnen zeggen, verheven tot zijn idealiteit.

Voor de spelers betekent dit dat zij hun rollen niet eenvoudigweg vervullen zoals in elk spel — in plaats daarvan spelen zij hun rollen, zij stellen deze voor voor het publiek. De manier waarop zij deelnemen aan het spel wordt niet langer bepaald door het feit dat zij volledig in het spel opgaan, maar door het feit dat zij hun rol spelen in relatie tot en met het oog op het geheel van het spel, waarin niet zij, maar het publiek moet opgaan. Er vindt een volledige verandering plaats wanneer spel als zodanig een toneelspel wordt. Het plaatst de toeschouwer op de plaats van de speler. Hij — en niet de speler — is de persoon voor wie en in wie het spel wordt gespeeld. Natuurlijk betekent dit niet dat de speler de betekenis van het geheel, waarin hij zijn voorstellende rol speelt, niet kan ervaren. De toeschouwer heeft slechts methodologische voorrang: doordat het spel voor hem wordt voorgesteld, wordt duidelijk dat het spel in zichzelf een te begrijpen betekenis draagt en daarom losgemaakt kan worden van het gedrag van de speler. In wezen is het verschil tussen speler en toeschouwer hier opgehoven. De eis dat het spel zelf in zijn betekenisvolheid bedoeld is, is voor beide hetzelfde.

Dit is nog steeds het geval, zelfs wanneer de speelgemeenschap zich afsluit voor alle toeschouwers, ofwel omdat zij zich verzet tegen de sociale institutionalisering van het artistieke leven, zoals bij de zogenaamde kamermuziek, die streeft naar authentieker musiceren door uitgevoerd te worden voor de spelers zelf en niet voor een publiek. Als iemand op deze manier muziek uitvoert, probeert hij in feite de muziek ``mooi te laten klinken'', maar dat betekent dat deze er in werkelijkheid voor elke luisteraar zou zijn. Artistieke presentatie bestaat van nature voor iemand, zelfs als er niemand is die louter luistert of toekijkt.

\subsection{Verandering tot structuur en totale bemiddeling}

\subsection{De temporaliteit van het esthetische}

\subsection{Het voorbeeld van de tragedie}

\section{Esthetische en hermeneutische consequenties}

\subsection{De ontologische betekenis van het beeld}

\subsection{De ontologische grondslag van het occasionele en het decoratieve}

\subsection{De grenspositie van de literatuur}

\subsection{Reconstructie en integratie als hermeneutische taken}
